%!TEX root = ../chapter04.tex 
%----------------------------------------------------------------------------------------
%	
%----------------------------------------------------------------------------------------

%----------------------------------------------------------------------------------------
%	 
%----------------------------------------------------------------------------------------


\newpage
\loadgeometry{marginlayout}  
\pagestyle{scrheadings} % Print headers again  
\KOMAoptions{
	headwidth=textwithmarginpar,
	headsepline=true,
	headsepline= 0.5pt,
	footwidth=textwithmarginpar,
}   
\onehalfspacing   




\section{Le raisonnement par l'absurde}

% \begin{tBox}[frametitle=Réfuter]
%	\textbf{Réfuter} une affirmation $P$  revient à démontrer que l'implication 
%	\begin{quote}
%		 \og $P$ implique FAUX \fg{}
%	\end{quote}
%\end{tBox}   
\begin{tBox}[frametitle=Démontrer par l'absurde] 
	\textbf{Démontrer par l'absurde}\sidenote{Raisonnement par l'absurde de Galilée sur la chute des corps. \href{https://youtu.be/dlDWHIX-c5M}{Vidéo d'Etiuenne Klein, France-Culture}.} que $P$ est VRAI revient à  démontrer l'implication :
	\begin{quote}
	 \og $(\operatorname{non}P)$ implique FAUX \fg{} 
	\end{quote} 
\end{tBox}

 
Par la suite nous utiliserons plusieurs fois le résultat suivant.
\begin{proposition}
	Pour tout entier $n\in\mathbb{N}$. $n$ et $n^2$ sont de même parité.
\end{proposition}
\begin{theorem}
	$\sqrt{2}$ est un irrationnel. 
	
	Il n'exite pas d'entiers $a, b\in\mathbb{N}$ tel que $2b^2=a^2$. 
\begin{marginfigure}[-2cm]
	\begin{tikzpicture}[scale=1]
		\tkzSetUpPoint[size=4, fill=black]
		\tkzfctset{tan style/.style={->,>=latex, color=red, line width=1.2pt}}   
		%			\tkzInit[xmin=-1.0,ymin=-1.0,xmax=10,ymax=6.0]
		%			\tkzClip
		\tkzDefPoint(0,0){A} 
		\tkzDefPoint(3,0){B}
		\tkzDrawTriangle[two angles= 45 and 90 ](A,B) \tkzGetPoint{C}
		\tkzLabelSegment[below](A,B){$b$}
		\tkzLabelSegment[right](B,C){$b$}
		\tkzLabelSegment[above left](A,C){$a$}
		\tkzMarkRightAngle[  size=0.3](C,B,A)  
		\tkzLabelPoint[below](A){\scriptsize$P$}
		\tkzLabelPoint[below right](B){\scriptsize$Q$}
		\tkzLabelPoint[above right](C){\scriptsize$R$} 
	\end{tikzpicture} 
	\caption{Il n'existe pas de triangle rectangle isocèle dont les côtés sont tous entiers.}
\end{marginfigure}
\end{theorem} 

\begin{proof}au programme.
	
Nous allons montrer \og $\sqrt{2}\in \mathbb{Q}$ implique  FAUX\fg{} en 3 étapes.

Supposons $\sqrt{2}\in\mathbb{Q}$. Il existe $a,b\in\mathbb{N}$   :
\begin{gather*}
	\sqrt{2}=\frac{a}{b} \qquad \text{$\frac{a}{b}$ irréductible}\\
	2 = \frac{a^2}{b^2} \\
	2 b^2 = a^2 
\end{gather*} 
\paragraph{Établir la parité de $a$ et de $b$}
$a^2$ est pair.\\
$a$ est de même parité que $a^2$, $a$ est aussi pair.\\
$b$ ne peut pas être pair, la fraction $\frac{a}{b}$ serait réductible.\\
$b$ est forcément impair.
\paragraph{Établir la contradiction à partir des parités de $a$ et $b$}
$a$ est pair : il existe $p\in\mathbb{N}$ tel que $a=2p$.\\
$b$ est impair : il existe $q\in\mathbb{N}$ tel que $b=2q+1$.\\ 
$\begin{WithArrows}
	2(2q+1)^2&=(2p)^2\Arrow{On développe}\\
	4q^2+4q+2&=4p^2\Arrow{$\times \frac{1}{2}$}\\
	2q^2+2q+1&=2p^2\\
	\underbrace{2(q^2+q)+1}_{\text{impair}}&=\underbrace{2p^2}_{\text{pair}}
	\end{WithArrows}$ 

D'ou la contradiction. FAUX.
\end{proof}

%----------------------------------------------------------------------------------------
%	
%----------------------------------------------------------------------------------------
\newpage
\loadgeometry{widelayout}  
\pagestyle{scrheadings} % Print headers again  
\KOMAoptions{
	headwidth=textwithmarginpar,
	headsepline=true,
	headsepline= 0.5pt,
	footwidth=textwithmarginpar,
}   
\onehalfspacing   

\newpage
\subsection{Exercices : raisonnement par l'absurde}
\begin{exercise}[$\sqrt{6}\notin\mathbb{Q}$]\footnote{La même méthode s'applique pour $\sqrt{10}$, $\sqrt{14}$ et $\sqrt{18}$.}   
	
Supposons le contraire : $\sqrt{6}$ est un rationnel et $\sqrt{6}=\frac{a}{b}$ irréductible avec $a, b\in\mathbb{N}$.
\begin{enumerate}[label={\arabic*)},leftmargin=*]
	\item Établir que $6b^2=a^2$
	\item 
	\begin{enumerate}[label={\alph*)},leftmargin=*] 
		\item Justifier soigneusement que $a$ est un nombre pair.
		\item En déduire que $b$ est un nombre impair.
	\end{enumerate}
	\item Soit $p$ et $q$ deux entiers tels que $a=2p$, et $b=2q+1$.
	\begin{enumerate}[label={\alph*)},leftmargin=*] 
		\item Montrer que $6q^2+6q+3=2p^2$
		\item Expliquer pourquoi la dernière égalité amène une contradiction qui permet de conclure.
	\end{enumerate}
\end{enumerate} 
\end{exercise}
\vfill 
\begin{exercise}[$\sqrt{3}\notin\mathbb{Q}$]\footnote{La même méthode s'applique pour $\sqrt{7}$, $\sqrt{11}$, $\sqrt{15}$ et $\sqrt{19}$.}\hfill 
	
Supposons le contraire : $\sqrt{3}$ est un rationnel et $\sqrt{3}=\frac{a}{b}$ irréductible avec $a, b\in\mathbb{N}$.
\begin{enumerate}[label={\arabic*)},leftmargin=*]
	\item Établir que $3b^2=a^2$
	\item
	\begin{enumerate}[label={\alph*)},leftmargin=*]  
		\item Justifier soigneusement que si $b$ est pair, alors $a$ est aussi pair.
		\item Justifier soigneusement que si $b$ est impair, alors $a$ est aussi impair.
		\item Expliquer pourquoi la possibilité que $a$ et $b$ soient tous les deux pairs est à rejeter.
		\item Que peut-on en déduire ?
	\end{enumerate}
	\item Soit $p$ et $q$ deux entiers tels que $a=2p+1$, et $b=2q+1$.
	\begin{enumerate}[label={\alph*)},leftmargin=*] 
		\item Montrer que $6q^2+6q+1=2p^2+2p$
		\item Expliquer pourquoi la dernière égalité amène une contradiction qui permet de conclure.
	\end{enumerate}
\end{enumerate} 
\end{exercise} 
\vfill 
\begin{exercise}[$\sqrt{5}\notin\mathbb{Q}$]\footnote{La même méthode s'applique pour $\sqrt{13}$, mais ne marchera pas pour $\sqrt{17}$ !}\hfill 
	
Supposons le contraire : $\sqrt{5}$ est un rationnel et $\sqrt{5}=\frac{a}{b}$ irréductible avec $a, b\in\mathbb{N}$.
\begin{enumerate}[label={\arabic*)},leftmargin=*]
	\item Établir que $5b^2=a^2$
	\item  En suivant la même démarche que la question 2 de l'exercice précédent, montrer que $a$ et $b$ sont tous les deux impairs.
	\item Soit $p$ et $q$ deux entiers tels que $a=2p+1$, et $b=2q+1$.
	\begin{enumerate}[label={\alph*)},leftmargin=*] 
		\item Montrer que $5q^2+q+1=p^2+p$, puis que $5q(q+1)+1=p(p+1)$.
		\item Justifier que $p(p+1)$ et $q(q+1)$ sont des nombres pairs.
		\item Expliquer pourquoi cela amène une contradiction qui permet de conclure.
	\end{enumerate}
\end{enumerate} 
\end{exercise} 
\vfill 
\begin{exercise}[$\sqrt{12}\notin\mathbb{Q}$]\footnote{La même méthode s'applique pour $\sqrt{8}=2\sqrt{2}$, $\sqrt{18}=3\sqrt{2}$ et $\sqrt{20}=2\sqrt{5}$.} \hfill 
	
Supposons le contraire : $\sqrt{12}$ est un rationnel et $\sqrt{12}=\frac{a}{b}$ irréductible avec $a, b\in\mathbb{N}$.
	\begin{enumerate}[label={\alph*)},leftmargin=*]
		\item Justifier que $\tfrac{\sqrt{12}}{2}$ est aussi un rationnel.
		\item En simplifiant la racine carré, montrer que $\tfrac{\sqrt{12}}{2}=\sqrt{3}$.
		\item Expliquer pourquoi cela amène une contradiction qui permet de conclure.
	\end{enumerate} 
\end{exercise} 
%\begin{remark}
%	$\sqrt{17}$ est un irrationnel, mais les approches précédentes ne permet pas de le prouver.
%\end{remark} 

%----------------------------------------------------------------------------------------
%	
%----------------------------------------------------------------------------------------
\newpage

\begin{exercise}\emoji{cactus}\hfill 
\begin{enumerate}[label={\alph*)},leftmargin=*]
	\item Expliquer pourquoi l'implication \og Si $x$ et $y$ sont rationnels alors $x+y$ est un rationnel \fg{} est vraie.
	\item Donner un contre-exemple pour réfuter l'implication  \og Si $x$ et $y$ sont irrationnels alors $x+y$ est un irrationnel \fg{}.
	\item L'argument suivant est-il valide ? 
	\begin{quote}
		$\sqrt{2}$ et $\sqrt{3}$ sont irrationnels, donc leur somme est aussi un irrationnel.
	\end{quote}
	\item Écrire la contraposée de l'implication de la question a. Est-elle vraie ?
	\item Démontrer par l'absurde que $\sqrt{2}+\sqrt{3}$ est un irrationnel.
\end{enumerate}
\end{exercise}



\begin{exercise}\hfill 
	
$n\in\mathbb{N}$ est un entier différent de 1. $x$ est le plus petit diviseur de $n$ autre que 1.  Nous voulons démontrer par l'absurde  que $x$ est toujours un nombre premier.
 
Recopier les encadrés en les mettant dans le bon ordre pour avoir une démonstration valide.

\noindent \tikz[baseline=(A.base)] \node[draw] (A) {Supposons au contraire que $x$ n'est pas premier}; \qquad
\tikz[baseline=(F.base)] \node[draw] (F) {$p$ et $q$ sont inferieur à $x$.}; \qquad
\tikz[baseline=(D.base)] \node[draw] (D) {$n=pqy$.}; \qquad
\tikz[baseline=(G.base)] \node[draw ] (G) {Ceci contredit l'hypothèse que $x$ est le plus petit diviseur différent de $1$ de $n$.}; \\
\tikz[baseline=(C.base)] \node[draw] (C) {il existe $y\in\mathbb{N}$ tel que $n=xy$}; \qquad
\tikz[baseline=(E.base)] \node[draw] (E) {$p$ et $q$ sont des diviseurs de $n$}; \\
\tikz[baseline=(B.base)] \node[draw] (B) {il existe $p$ et $q\in\mathbb{N}$ différents de 1 tel que $d=pq$.};   
\end{exercise}
\begin{example}[Application] En trouvant le plus petit facteur autre que 1 d'un nombre, on trouve aussi le plus petit facteur premier. Ce théorème est à la base des algorithmes de décomposition en facteurs premiers d'un entier. 
	
Analysons le fonctionnement des instructions \lstinline|decomposition(24)| et \lstinline|decomposition(150)| :
 
\noindent\begin{lstlisting}[ language=python , 
	%	title=Programme A,   
	linewidth=0.95\columnwidth,
	style=python-code,    
	aboveskip=0.2\topskip,
	belowskip=0.2\topskip,
%	lineskip=1.1em,
	xleftmargin=1.2em,
	tabsize=4,
	%	firstnumber=1  
	]
def decomposition(n) :		# n est un entier
	x = 2 
	while n != 1 :			#	
		if n%x==0 :			# 
			print(x)		#
			n = n // x		#
		else :
			x = x + 1		#
	return "FIN"
\end{lstlisting}
 
\medskip
\noindent\Decomposition[TableauVerticalVide]{24}	\hfill 
\Decomposition[TableauVerticalVide]{150} \hfill\null 

\medskip
\noindent Déterminer le plus grand facteur premier de \numprint{600851475143}\dotfill %  \numprint{101010101010101010101} 
\end{example}
 
%----------------------------------------------------------------------------------------
%	 
%----------------------------------------------------------------------------------------

\cleardoublepage